\section{研究对象与方法}
\label{sec:method}

\subsection{示例程序与语义覆盖}

研究对象应是一个规模小但覆盖面明确的 SysY 程序。计划采用“数组聚合后分支修正”的单一程序：辅助函数通过循环遍历数组并累加，主函数读入数据、调用辅助函数，再依据阈值选择输出路径。表~\ref{tab:coverage}给出特征到证据的映射。最终源文件路径和确定输入应在实验实现后补入，而非在此稿中臆造。

\begin{table}
\caption{语言特征、低层结构与验证证据的对应关系}
\label{tab:coverage}
\centering
\begin{tabular}{lll}
\toprule
SysY 特征 & LLVM IR / RISC-V 观察点 & 验证证据 \\
\midrule
数值运算与赋值 & 算术指令、寄存器或栈槽 & 已保存 IR/汇编 \\
条件分支 & 基本块、条件跳转 & 控制流图与路径用例 \\
循环 & 回边、归纳变量 & 边界与多次迭代用例 \\
函数调用 & \texttt{call}、调用约定 & 参数/返回值一致性 \\
数组访问 & \texttt{getelementptr}、地址计算 & 多元素输入用例 \\
运行时输入输出 & 外部符号与链接重定位 & 标准输出和链接映射 \\
\bottomrule
\end{tabular}
\end{table}

\subsection{四阶段工具链观察}

对同一源文件分别停止在预处理、编译、汇编和链接之后，保留以下制品：
\begin{description}
  \item[预处理器（preprocessor）] 展开头文件、宏与条件编译，产生仍接近源语言的翻译单元；通过预处理输出确认声明的来源。
  \item[编译器（compiler）] 完成词法、语法与语义分析，构造中间表示并执行优化，最终生成目标汇编；通过未优化和优化 IR 的差分观察结构变化。
  \item[汇编器（assembler）] 将助记符编码为可重定位目标文件，生成符号表和重定位项；通过反汇编检查机器指令与源汇编的对应。
  \item[链接器（linker）] 解析符号、应用重定位并组合运行时库，产生可执行文件；通过符号表和链接映射确认 SysY 运行时接口的绑定。
\end{description}

这种分层描述刻意区分工具边界与编译器内部阶段。ELF（Executable and Linkable Format）目标文件中的节、符号和重定位项，是汇编器与链接器之间可直接检查的契约\cite{elfspec2024}。

\subsection{三种实现与差分判定}

以 SysY/参考编译路径作为基线，人工编写等价 LLVM IR，并以 RISC-V 指令集架构（Instruction Set Architecture, ISA）实现汇编版本\cite{riscv20240411}。LLVM IR 使用不透明指针（opaque pointer）语法，以匹配当前 LLVM 版本的表示约定\cite{llvmLangRef}; RISC-V 版本遵守工具链选定的应用二进制接口（Application Binary Interface, ABI）。

对每个测试夹具，驱动脚本记录标准输入、标准输出、标准错误和退出状态。判定函数为
\begin{equation}
  \operatorname{pass}(i)=
  \bigwedge_{p \in \{\mathrm{source},\mathrm{IR},\mathrm{asm}\}}
  \left(O_{p,i}=O_{\mathrm{oracle},i}\ \land\ E_{p,i}=E_{\mathrm{oracle},i}\right),
\end{equation}
其中 $O$ 是规范化后的标准输出，$E$ 是退出状态。规范化仅允许消除约定明确的行尾差异，不能吞掉空白、诊断或数值差异。测试集至少覆盖零次循环、一次循环、多次循环、条件两侧路径及数组边界附近的合法输入。

\subsection{优化对照}

在源程序、目标三元组（target triple）和运行时库不变时，对比 \texttt{-O0} 与一个已记录的优化级别。比较维度包括基本块数、静态指令数、栈访问和关键优化现象；性能只在有重复次数、预热策略与计时方法时报告。优化后汇编更短并不自动意味着更快，因而结构指标与运行时间必须分开解释。
